\section{宇宙}\etiquette[0]{I.sec0}

\emph{宇宙（universe）}是满足下列性质的非空集合\nde{这一定义与\textit{下文} Bourbaki 的附录相矛盾. 不过, 空宇宙是否有用似乎值得商榷……} $𝒰$:
\begin{itemize}
\item[($𝒰$ 1)] 若 $x\in 𝒰$ 且 $y\in x$, 则 $y\in 𝒰$.

\item[($𝒰$ 2)] 若 $x$, $y\in 𝒰$, 则 $\{x,y\}\in 𝒰$.

\item[($𝒰$ 3)] 若 $x\in 𝒰$, 则 $\Pp(x)\in 𝒰$.

\item[($𝒰$ 4)] 若 $(x_i, i\in I\sisi{\in 𝒰}{})$ 是 $𝒰$ 的一族元素\sisi{}{且 $I\in𝒰$}, 则 $\cup_{i\in I}x_i \in
  𝒰$.
\end{itemize}

由\pageoriginale{}上述公理容易推出下列性质:
\begin{itemize}
\item 若 $x\in 𝒰$, 则集合 $\{x\}$ 属于 $𝒰$.

\item 若 $x$ 是 $y\in 𝒰$ 的子集, 则 $x\in 𝒰$.

\item 若 $x$, $y\in 𝒰$, 则有序对 $(x,y)=\{\{x\},\{x,y\}\}$ (Kuratowski 的定义) 是 $𝒰$ 的元素.

\item 若 $x$, $y\in 𝒰$, 则并集 $x\cup y$ 与 %%
积 $x\times
   y$ 都是 $𝒰$ 的元素.

\item 若 $(x_i, i\in I\in 𝒰)$ 是 $𝒰$ 的一族元素,
则积 $\prod_{i\in I}x_i$ 是 $𝒰$ 的元素.

\item 若 $x\in 𝒰$, 则 $\card(x)< \card(𝒰)$ (严格不等式). 特别地, 关系 $𝒰\in 𝒰$ 不成立.
\end{itemize}

% END INCLUDED SGA4_English_current_buildable_source_20260906/source/01/drafts/SGA4_I_LINES_46_100_English.texfrag


% BEGIN INCLUDED SGA4_English_current_buildable_source_20260906/source/01/drafts/SGA4_I_LINES_102_196_English.texfrag
因此, 从一个宇宙的元素出发, 可以进行集合论中所有通常的运算,
而最终结果仍然是这个宇宙的元素.

宇宙概念的第一个好处是, 它给出了通常所用范畴的定义:
属于宇宙 $𝒰$ 的集合所成的范畴 ($𝒰$-$\Ens$),
属于宇宙 $𝒰$ 的拓扑空间所成的范畴,
属于宇宙 $𝒰$ 的交换群所成的范畴
($𝒰$-$\Ab$), 属于宇宙 $𝒰$ 的范畴所成的范畴\ldots.

然而, 唯一已知的宇宙是形如
$\{\{∅\}, \{\{∅\}, ∅\}\}$ \etc 的符号所成的集合
(这个宇宙的所有元素都是有限集, 而且这个宇宙是可数的).
特别地, 尚不知道任何含有无限基数元素的宇宙.
因此, 我们在集合论公理中加入公理:
$(𝒰 A)$ \emph{对每个集合 $x$, 都存在宇宙 $𝒰$ 使得 $x\in
  𝒰$.}

由于一族宇宙的交\pageoriginale{}仍是宇宙, 立即可知,
每个集合都是某个最小宇宙的元素.
可以证明, 公理 $(𝒰
A)$ 独立于集合论公理.

我们还要加入公理:


\noindent \emph{$(𝒰 B)$ 设 $R\{x\}$ 是一个关系, $𝒰$ 是一个宇宙.
若存在元素 $y\in 𝒰$ 使得 $R\{y\}$, 则
  $\tau_xR\{x\}\in 𝒰$.}

公理 $(𝒰 A)$ 和 $(𝒰 B)$ 相对于集合论其余公理的相容性尚未得到证明,
而且看来也不可能得到证明.

设 $𝒰$ 是一个宇宙, $c(𝒰)$ 是 $𝒰$ 的元素的基数的上确界
($c(𝒰)\leq\card(𝒰)$). 基数 $c(𝒰)$ 有下列性质:

$(FI)$ 若 $a<c(𝒰)$, 则 $2^a<c(𝒰)$.

$(FII)$ 若 $(a_i, i\in I)$ 是一族严格小于 $c(𝒰)$ 的基数,
且 $\card(I)$ %%
严格小于 $c(𝒰)$, 则 $\Sigma_{i\in I} a_i <
c(𝒰)$.

具有性质 $(FI)$ 和 $(FII)$ 的基数称为\emph{强不可达基数（strongly inaccessible cardinal）.}

基数 0 和可数无限基数都是强不可达的.

公理 $(𝒰 A)$ 蕴含:

\emph{$(𝒰 A')$ 每个基数都严格小于某个强不可达基数.}

反过来, 可以证明 (11), $(UA')$ 的相容性蕴含 $(UA)$ 的相容性,
而这些公理的相容性又蕴含公理 $(𝒰 B)$ 的相容性.

如果\pageoriginale{}一个集合 $E$ 不存在满足 $x_0 \in E$, $x_{n+1}\in x_n$ 的无限族
$(x_n, n\in N)$, 就称它为\emph{Artin 集（Artinian set）}.
可以证明 [\loccit], 强不可达基数与 Artin 宇宙之间有如下的一一对应:
给每个强不可达基数 $c$ 对应唯一的 Artin 宇宙 $U_c$, 满足
$$
\card(𝒰_c)=c{\quoi}.
$$
设 $c<c'$ 是两个强不可达基数. 我们有:
$$
𝒰_c\in 𝒰_{c'}{\quoi}.
$$
特别地, 基数小于一个给定基数的 Artin 宇宙构成一个良序集
(关于属于关系而言). 公理 $(𝒰 A')$ 等价于公理:

$(𝒰 A'')$. \emph{每个 Artin 集都是某个 Artin 宇宙的元素.}

注意, 所有通常所用的集合 $(\ZZ, \QQ,\ldots)$ 都是 Artin 集.

% END INCLUDED SGA4_English_current_buildable_source_20260906/source/01/drafts/SGA4_I_LINES_102_196_English.texfrag


% BEGIN INCLUDED SGA4_English_current_buildable_source_20260906/source/01/drafts/SGA4_I_LINES_197_410_English.texfrag
\section{$𝒰$-范畴. 集合预层}\etiquette[1]{I.sec1}

\setcounter{subsection}{-1}
\subsection{}\etiquette[1.0]{I.sec1.0}
\emph{在本讨论班的其余部分, 除非另有明确说明,
所考虑的宇宙都含有一个具有无限基数的元素.}
设 $𝒰$ 是一个宇宙. 如果一个集合\emph{同构于} $𝒰$ 的一个元素,
就称它是 $𝒰$-\emph{小的（small）}
(或者在不致混淆时, 简称\emph{小的}). 我们也将使用小群、小环、
小范畴\nde{将一个范畴 $C$ 看作箭头的集合
(配备对象子集, 将其视为恒等箭头的集合, 并配备“源、靶”映射).
因此, “$C$ 是 $𝒰$ 的元素”以及“$C$ 是 $𝒰$-小的”这些说法都有意义.}
\ldots 这样的术语. 我们常常不加明说地假定,
所处理的概形、拓扑空间、指标集\ldots 都 $\in 𝒰$,
或\pageoriginale{}至少它们的基数 $\in 𝒰$;
不过, 我们要处理的许多范畴并不 $\in 𝒰$.

\begin{enonce*}{定义 1.1}\etiquette[1.1]{I.def1.1}
设 $𝒰$ 是一个宇宙, $C$ 是一个范畴.
若对 $C$ 的每对对象 $(x,y)$, 集合 $\Hom_C(x,y)$ 都是 $𝒰$-小的,
就称 $C$ 是一个 $𝒰$-范畴.
\end{enonce*}

\setcounter{subsection}{1}
\subsubsection{}\etiquette[1.1.1]{I.sec1.1.1}
设 $C$ 和 $D$ 是两个范畴, $\text{Fonct}(C,D)$ 是从 $C$ 到 $D$ 的函子所成的范畴.
下列断言是显然的:
\begin{enumerate}
\item[a)] 若 $C$ 和 $D$ 是宇宙 $𝒰$ 的元素
(\resp 是 $𝒰$-小的), 则范畴 $\text{Fonct}(C,D)$ 是
$𝒰$ 的元素 (\resp 是 $𝒰$-小的).

\item[b)] 若 $C$ 是 $𝒰$-小的, 且 $D$ 是一个 $𝒰$-范畴,
则 $\text{Fonct}(C,D)$ 是一个 $𝒰$-范畴.
\end{enumerate}

\begin{enonce*}[remark]{注 1.1.2}\etiquette[1.1.2]{I.rem1.1.2}
设 $D$ 是具有下列性质的范畴:
\begin{itemize}
\item[(C1)] 集合 $\ob(D)$ 包含于宇宙 $𝒰$.

\item[(C2)] 对 $D$ 的每对对象 $(x,y)$,
集合 $\Hom_D(x,y)$ 是 $𝒰$ 的元素.
\end{itemize}
(从宇宙 $𝒰$ 构造的通常所用范畴都有这两个性质:
$𝒰$-$\Ens$,
$𝒰$-$\Ab$,\ldots). 设 $C$ 是属于 $𝒰$ 的范畴.
则范畴 $\text{Fonct}(C,D)$ 一般不具有性质 (C1) 和 (C2).
例如, 范畴 $\text{Fonct}(C,𝒰$-$\Ens$) 既不具有性质
(C1), 也不具有性质 (C2). 因此, 我们采用上述 $𝒰$-范畴的定义,
而不采用由上面条件 (C1) 和 (C2) 给出的限制更强的概念.
\end{enonce*}

\begin{enonce*}{定义 1.2}\etiquette[1.2]{I.def1.2}
设\pageoriginale $C$ 是一个范畴.
$C$ 上相对于宇宙 $𝒰$ 的集合预层（presheaf of sets）所成的范畴
(或者在不致混淆时, 简称 $C$ 上的预层范畴)
是 $C$ 上取值于 $𝒰$-集合范畴的反变函子所成的范畴.
\end{enonce*}

我们用 $\widehat{C}_{𝒰}$
(或者在不致混淆时, 更简单地用 $\widehat{C}$\,)
表示 $C$ 上相对于宇宙 $𝒰$ 的集合预层所成的范畴.
$\widehat{C}_{𝒰}$ 的对象称为 $C$ 上的 $𝒰$-预层
(或简称预层). 当 $C$ 是 $𝒰$-小的时,
范畴 $\widehat{C}_{𝒰}$ 是一个 $𝒰$-范畴.
当 $C$ 是一个 $𝒰$-范畴时, $\widehat{C}_{𝒰}$ 未必是 $𝒰$-范畴.


\begin{enonce*}[remark]{构造--定义
    1.3}\etiquette[1.3]{I.condef1.3}
设 $x$ 是 $𝒰$-范畴 $C$ 的一个对象.
\emph{由 $x$ 表示的 $𝒰$-函子}是如下构造的函子
$h_{𝒰}(x):C^\circ\rightarrow 𝒰$-$\Ens$
\footnote[1]{$C^\circ$ 总表示 $C$ 的反范畴.}.
设 $y$ 是 $C$ 的一个对象.
\begin{enumerate}
\item[a)] 若 $\Hom_C(y,x)$ 是 $𝒰$ 的元素, 则:
      $$
      h_{𝒰}(x)(y)=\Hom_C(y,x)
      $$

\item[b)] 假定 $\Hom_C(y,x)$ 不是 $𝒰$ 的元素,
并以 $R(Z,x,y)$ 表示关系:
“集合 $Z$ 是某个同构 $\Hom_C(y,x)\xrightarrow{\sim}Z$ 的终域”. 于是置:
  $$
  h_{𝒰}(x)(y)=\tau_ZR(Z){\quoi}.
  $$
\end{enumerate}
(由公理 $(𝒰 B)$, $h_{𝒰}(x)(y)$ 是 $𝒰$ 的元素).
以 $R'(u,x,y)$ 表示关系: “$u$ 是双射
$$
u:\Hom_C(y,x)\xrightarrow{\sim} h_{𝒰}(x)(y)\mbox{
”}{\quoi}.
$$
于是\pageoriginale{}置:
$$
\varphi(y,x)=\tau_uR'(u){\quoi}.
$$

注意, 在情形 (a) 和 (b) 中, 都有一个典范同构:
$$
\varphi(y,x):\Hom_C(y,x)\xrightarrow{\sim}
h_{𝒰}(x)(y){\quoi}.
$$
(在情形 a) 中, $\varphi(y,x)$ 是恒等映射).
再设 $u:y\rightarrow y'$ 是 $C$ 的一个箭头.
由态射的复合, 态射 $u$ 定义映射:
$$
\Hom_C(u,x):\Hom_C(y',x)\longrightarrow \Hom_C(y,x){\quoi}.
$$
于是置
$$
h_{𝒰}(x)(u)=\varphi(y,x)\Hom_C(u,x)\varphi(y',x)^{{}-1}{\quoi}.
$$

显然, 这样定义的 $h_{𝒰}(x)$ 是函子
$C^\circ\rightarrow 𝒰$-$\Ens$.
\end{enonce*}

\setcounter{subsection}{3} \setcounter{subsubsection}{0}
\subsubsection{}\etiquette[1.3.1]{I.sec1.3.1}
同样, 利用同构 $\varphi$, 对每个态射 $y:x\rightarrow x'$,
定义函子之间的态射:
$$
h_{𝒰}(y):h_{𝒰}(x)\longrightarrow h_{𝒰}(x')
$$
显然, 这就定义了一个函子
$$
h_{𝒰}:C\longrightarrow
\widehat{C}_{𝒰}{\quoi}.
$$

\subsubsection{}\etiquette[1.3.2]{I.sec1.3.2}
再设 $𝒱$ 是一个宇宙, 且 $𝒰\subset
𝒱$. 于是有一个典范的全忠实函子:
$$
𝒰\text{-}\Ens\hookrightarrow 𝒱\text{-}\Ens
$$
从而有全忠实函子:
$$
\widehat{C}_{𝒰}\hookrightarrow
\widehat{C}_{𝒱}
$$
而\pageoriginale{}图:

%\def\objectstyle{\mathrm}--here
\[
\xymatrix@R=1.75cm@C=3.5cm@M=.2cm{
C\ar[r]^{h_{\underline{U}}} \ar[dr]_{h_{\underline{V}}}& C{\widehat{\phantom{i}}}_{\underline{U}}\ar@{_{(}->}[d]\\
& C\widehat{\phantom{i}}_{\underline{V}} }
\]
在典范同构意义下交换. 当 $C$ 是 $𝒰$ 的元素时,
这个典范同构是恒等同构.

\subsubsection{}\etiquette[1.3.3]{I.sec1.3.3}
\emph{实际使用时, 宇宙 $𝒰$ 一经选定就固定下来, 不再提及.}
于是我们采用记号 $\widehat{C}$ (表示 $𝒰$-集合预层的范畴) 以及
$$
h:C\longrightarrow \widehat{C}{\quoi}.
$$
对 $C$ 的每个对象 $x$, 预层 $h(x)$ 称为\emph{由 $x$ 表示的预层},
并且我们总将预层 $h(x)$ 在 $y\in \ob(C)$ 处的值等同于 $\Hom_C(y,x)$.

\begin{enonce*}{命题 1.4}\etiquette[1.4]{I.prop1.4}
设 $C$ 是一个 $𝒰$-范畴, $F$ 是 $C$ 上的预层,
$X$ 是 $C$ 的一个对象. 有一个关于 $X$ 和 $F$ 均函子性的同构:
$$
i:\Hom_{\widehat{C}}(h(X),F)\xleftarrow{\sim}
F(X){\quoi},
$$
其中, 当 $F$ 具有 $h(Y)$ 的形式时, 映射 $i$ 正是映射:
$$
h:\Hom(h(X), h(Y))\longleftarrow \Hom(X,Y){\quoi}.
$$
特别地, 函子 $h$ 是全忠实的.
\end{enonce*}

\setcounter{subsection}{4} \setcounter{subsubsection}{0}
\subsubsection{}\etiquette[1.4.1]{I.sec1.4.1}
这一命题说明, 将 $C$ 的对象与相应的反变函子等同的惯常说法是合理的.
与 $h$ 的像中的一个对象同构
(或者沿用上述惯常说法, 与 $C$ 的一个对象同构) 的预层\pageoriginale,
称为\emph{可表示预层（representable presheaf）.}

\subsubsection{}\etiquette[1.4.2]{I.sec1.4.2}
设 $𝒱$ 是包含宇宙 $𝒰$ 的宇宙.
$𝒰$-预层范畴是 $𝒱$-预层范畴的一个\emph{满}子范畴.
因此, 一个 $𝒰$-预层可表示, 当且仅当它在 $𝒱$-预层范畴中的像是可表示的 $𝒱$-预层.

% END INCLUDED SGA4_English_current_buildable_source_20260906/source/01/drafts/SGA4_I_LINES_197_410_English.texfrag


% BEGIN INCLUDED SGA4_English_current_buildable_source_20260906/source/01/drafts/SGA4_I_LINES_411_824_English.texfrag
\section{投射极限与归纳极限}\etiquette[2]{I.sec2}

设 $C$ 是一个 $𝒰$-范畴, $I$ 是一个小范畴.
以 $\Fonct(I,C)$ 表示从 $I$ 到 $C$ 的函子所成的范畴.
范畴 $\Fonct(I,C)$ 是一个 $𝒰$-范畴.
给 $C$ 的每个对象 $X$ 对应 $C$ 的子范畴 $𝒳$,
其唯一对象是 $X$, 唯一箭头是 $X$ 的恒等箭头.
以 $i_X:𝒳\rightarrow C$ 表示包含函子.
有唯一的函子 $e_X:I\rightarrow 𝒳$,
我们用 $k_X:I\rightarrow C$ 表示函子 $i_X\circ
e_X$. (称 $k_X$ 是取值为 $X$ 的常值函子.)
对应 $X\mapsto k_X$ 显然关于 $X$ 是函子性的,
因而对每个函子 $G:I\rightarrow
C$, 可以定义预层 $X\mapsto \Hom_{\Fonct(I,C)}(k_X,G)$.

\begin{enonce*}{定义 2.1}\etiquette[2.1]{I.def2.1}
$G$ 的投射极限（projective limit）, 记为 $\varprojlim_I G$, 是预层:
$$
X\longmapsto \Hom_{\Fonct(I,C)}(k_X,G){\quoi}.
$$
当预层 $\varprojlim_I G$ 可表示时,
我们也用 $\varprojlim_I G$ 表示 $C$ 中表示它的一个\pageoriginale{}对象.
因此, 对象 $\varprojlim_I G$ 仅在同构意义下确定.
在不致混淆时, 使用简写记号 $\varprojlim G$.
\end{enonce*}

当 $G$ 变化时, $\varprojlim G$ 是从范畴
$\Fonct(I,C)$ 到 $\widehat{C}$ 的函子.

\setcounter{subsection}{1} \setcounter{subsubsection}{0}
\subsubsection{}\etiquette[2.1.1]{I.sec2.1.1}
由对称性 (将 $C$ 中箭头的方向反转),
类似地定义函子的归纳极限（inductive limit）:
它是 $C$ 上取值于 $𝒰$-集合范畴的协变函子.
我们将使用记号 $\varinjlim_I
G$ 或 $\varinjlim G$.

注意, 积、纤维积和核都是投射极限.
同样, 和、融合和以及余核都是
%%
归纳极限.

\begin{enonce*}{定义 2.2}\etiquette[2.2]{I.def2.2}
设 $I$ 和 $C$ 是两个范畴, $G:I\rightarrow C$ 是函子.
如果存在宇宙 $𝒰$ 使得:
\begin{enumerate}
\renewcommand{\labelenumi}{\theenumi)}
\item 范畴 $I$ 是 $𝒰$-小的.

\item 范畴 $C$ 是一个 $𝒰$-范畴.

\item 取值于 $𝒰$-集合范畴的预层 $\varprojlim G$ 可表示.
\end{enumerate}
就称 $G$ 的投射极限可表示.
\end{enonce*}

由第 ${\rm n^\circ}$ 1 节可知,
表示预层 $\varprojlim G$ 的对象 $\varprojlim G$
在同构意义下与宇宙 $𝒰$ 无关.
还要注意, 存在具有性质 (1) 和 (2) 的最小宇宙 $𝒰'$,
而且预层 $\varprojlim G$ 必然取值于 $𝒰'$-集合范畴.

\setcounter{subsection}{2} \setcounter{subsubsection}{0}
\subsubsection{}\etiquette[2.2.1]{I.sec2.2.1}
设\pageoriginale $I$ 和 $C$ 是两个范畴.
如果对每个函子 $G:I\rightarrow C$, $G$ 的投射极限都可表示,
就说\emph{$C$ 中的 $I$-投射极限可表示}.
最后, 设 $C$ 是一个范畴, $𝒰$ 是一个宇宙.
如果对每个 $𝒰$-小范畴 $I$ 以及每个函子 $G:I\rightarrow
C$, $G$ 的投射极限都可表示,
就说\emph{$C$ 中的 $𝒰$-投射极限可表示}.

\begin{enonce*}{命题 2.3}\etiquette[2.3]{I.prop2.3}
设 $C$ 是一个范畴, $𝒰$ 是一个\sisi{非空的}{}宇宙.
下列断言等价:
\begin{enumerate}
\renewcommand{\theenumi}{\roman{enumi}}
\renewcommand{\labelenumi}{\theenumi)}
\item $C$ 中的 $𝒰$-投射极限可表示.

\item 以小集为指标的积可表示, 并且一对箭头的核可表示.
\item 以小集为指标的积可表示, 并且纤维积可表示.
\end{enumerate}
\end{enonce*}

\begin{enonce*}[remark]{证明}
只须注意到有一个关于 $G$ 函子性的同构
$$
\varprojlim G=\Ker (\prod_{i\in \ob(I)} %%
G(i)\rightrightarrows \prod_{u\in \rm{F\ell}(I)} G(\but(u))){\quoi},
$$
其中这对箭头由下列态射定义
\begin{align*}
\prod_{i\in \Ob(I)}  G(i) %%
 & \xrightarrow{\pr_{\but(u)}}{} G(\but(u))\\
\prod_{i\in \Ob(I)} G(i) %%
 &\xrightarrow{\pr_{\text{source}(u)}}{} G(\text{source
}(u))\xrightarrow{G(u)} G(\but(u)){\quoi}.
\end{align*}
此外, 显然 $\Ker (X  \displaystyle \Lright^{u}_{v}Y)=
X\prod_{X\times Y} X$, 其中从 $X$ 到 $X\times Y$ 的两个态射为
$\id_X \times u$ 和 $\id_X\times v$.
\end{enonce*}

\setcounter{subsection}{3} \setcounter{subsubsection}{0}
\subsubsection{}\etiquette[2.3.1]{I.sec2.3.1}
当然, 对归纳极限也有类似的定义和断言, 这里不再详述.
有限投射极限\pageoriginale{}和有限归纳极限
(\ie 关于有限范畴 $I$ 的极限) 也是如此.


\begin{enonce*}{推论 2.3.2}\etiquette[2.3.2]{I.cor2.3.2}
以 $𝒰$-$\Ens$ 表示 $𝒰$-集合的范畴,
以 $𝒰$-$\Ab$ 表示 $𝒰$-$\Ens$ 中的 Abel 群对象所成的范畴.
$𝒰$-$\Ens$ 和 $𝒰$-$\Ab$ 中的 $𝒰$-投射极限与归纳极限均可表示.
\end{enonce*}

\begin{enonce*}{命题 2.3.3}\etiquette[2.3.3]{I.prop2.3.3}
设 $I$ 是一个小范畴, $F:I\rightarrow 𝒰$-$\Ens$
是函子, 且 \linebreak %%
$G\subset\ob(I)$ 是 $I$ 的一个小对象集,
使得对 $I$ 的每个对象 $X$, 都存在对象 $Y\in G$ 及态射
$Y\rightarrow X$ (\resp $X\rightarrow Y$).
则 $\varprojlim_I F$
(\resp $\varinjlim_I F$) 可表示, 并且
$\card(\varprojlim_I F)\leq \Pi_{Y\in G}\card(F(Y))$ (\resp
$\card(\varinjlim_I F)\leq \Sigma_{Y\in G}\card(F(Y))$).
\end{enonce*}

\begin{enonce*}[remark]{证明}
立即可验证, $\varprojlim_I F$ (\resp $\varinjlim_I
F$) 同构于 $\Pi_{Y\in
G}F(Y)$ (\resp $\amalg_{Y\in G}F(Y)$) 的一个子对象
(\resp 一个商); 因而得到命题.
\end{enonce*}

\begin{enonce*}{定义 2.4.1}\etiquette[2.4.1]{I.def2.4.1}
设 $C$ 是一个有限投射极限 (\resp 归纳极限) 可表示的范畴,
$u:C\rightarrow C'$ 是函子.
若它与有限投射极限 (\resp 归纳极限) “交换”,
就称 $u$ 是左正合的 (\resp 右正合的).
既左正合又右正合的函子称为正合函子.
\end{enonce*}

\setcounter{subsection}{4} \setcounter{subsubsection}{1}
\subsubsection{}\etiquette[2.4.2]{I.sec2.4.2}
由 \hyperref[I.prop2.3]{2.3} 可知, 一个函子\sisi{}{左}正合, %% étiquette !
当且仅当它将终对象 (= 空积) 映为终对象,
将两个对象的积映为它们的像的积,
并将一对箭头的核映为像箭头对的核;
等价地, 它将终对象映为终对象, 并将纤维积映为纤维积\sisi{. (}{ (}这里假定 $C$ 中有限 $\varprojlim$ 都可表示\sisi{.)}{).}

\setcounter{subsection}{5} \setcounter{subsubsection}{-1}
\subsubsection{}\etiquette[2.5.0]{I.sec2.5.0}
设\pageoriginale $I$, $J$ 和 $C$ 是三个范畴,
$G:I\times J\rightarrow C$ 是一个函子
(\ie 从 $I$ 到从 $J$ 到 $C$ 的函子范畴的一个函子).
假定函子
\begin{align*}
G_i:&J\longrightarrow C\qquad i\in \ob(I)\\
    & j\longmapsto G(i\times j)
\end{align*}
的投射极限可表示, 并且函子
$$
i\longmapsto \varprojlim G_i
$$
具有可表示的投射极限.
则显然, 函子 $G$ 具有可表示的投射极限, 并有典范同构
$$
\varprojlim_{I\times J} G\\
\xrightarrow{\sim}\varprojlim_I \varprojlim_J G_i{\quoi},
$$
从而有典范同构
$$
\varprojlim_I \varprojlim_J G_i\xrightarrow{\sim}\varprojlim_J
\varprojlim_I G_j{\quoi}.
$$
下文中, 我们将简洁地说投射极限与投射极限交换.
同样, 归纳极限与归纳极限交换.
但一般说来, 归纳极限与投射极限并不交换.

\setcounter{subsection}{5}
\begin{enonce*}{定义 2.5}\etiquette[2.5]{I.def2.5}
设 $C$ 是一个纤维积可表示的范畴, $I$ 是一个范畴,
$G:I\rightarrow C$ 是函子,
$g:G\rightarrow X$ 是从 $G$ 到 $C$ 的一个对象的态射
(\ie 从 $G$ 到与 $X$ 对应的常值函子的态射),
$m:Y\rightarrow X$ 是 $C$ 中的一个态射.
设\pageoriginale $\fprod{G}{Y}{X}:I\rightarrow C$ 是函子
$$
i\longmapsto \fprod{G(i)}{Y}{X} {\quoi}.
$$
如果对每个对象 $X$、每个态射 $g:G\rightarrow X$
以及每个态射 $m:Y\rightarrow X$,
\begin{enumerate}
\item[a)] 函子 $\fprod{G}{Y}{X}$ 的归纳极限可表示,

\item[b)] 典范态射 $\varinjlim (\fprod{G}{Y}{X}
  )\rightarrow \fprod{(\varinjlim G)}{Y}{X}$ 是同构,
\end{enumerate}
就称 $G$ 的归纳极限是泛的（universal）.
\end{enonce*}

\begin{enonce*}{命题 2.6}\etiquette[2.6]{I.prop2.6}
设 $𝒰$ 是一个宇宙.
$𝒰$-$\Ens$ 中可表示的归纳极限
(特别地, $𝒰$-$\Ens$ 中的 $𝒰$-归纳极限 $\hyperref[I.sec2.2.1]{(2.2.1)}$) %% étiquette !
都是泛的.
\end{enonce*}
我们还将使用投射极限与归纳极限的另一个交换结果, 现在将其陈述如下.

\begin{enonce*}{定义 2.7}\etiquette[2.7]{I.def2.7}
若范畴 $I$ 具有下列性质, 就称它是伪滤过的（pseudo-filtered）:
\begin{itemize}
\item[PS 1)] 每个形如
\[
\xymatrix@R=.5cm@C=1cm{
       & j\\
i \ar[ru]\ar[dr]     &  \\
       & j'
}
\]
的图都能嵌入交换图
\[
\xymatrix@R=.5cm@C=1cm{
  & j \ar[rd]&\\
i\ar[ur]\ar[rd] &   & k\\
  & j'\ar[ru]&
}
\]

\item[PS 2)] 每个形如
\[
\xymatrix@R=.75cm@C=1.5cm@M=.25cm{
i\ar@<2.5pt>[r]^{u}\ar@<-2.5pt>[r]_{v} & j }
\]
%%
的图都能\pageoriginale{}嵌入一个图
\[
\xymatrix@C=1.75cm@M=.25cm{ i\ar@<2.5pt>[r]^{u}\ar@<-2.5pt>[r]_{v} &
j\ar[r]%%
^{w} & k\\
}
\]
使得
  $$
  w\circ u=w\circ v{\quoi}.
  $$
\end{itemize}
如果范畴 $I$ 伪滤过、非空且连通, 就称它是滤过的（filtered）.
这里连通是指 $I$ 的任意两个对象都可由一个箭头序列连接
(对箭头的方向不作任何限制);
在具有 {\rm PS 2} 的条件下,
这也就是说 $I\neq ∅$, 且对 $I$ 的任意两个对象 $a$, $b$,
总存在 $I$ 的一个对象 $c$ 及箭头 $a\sisi{$
$}{\rightarrow}c$ 和 $b\sisi{$ $}{\rightarrow}c$.
如果 $I^\circ$ 是滤过的, 我们也说范畴 $I$ 是余滤过的.
\end{enonce*}

\begin{enonce*}{例 2.7.1}\etiquette[2.7.1]{I.exem2.7.1}
如果 $I$ 中的融合和 (\resp 两个对象的和)
以及平行箭头对的余核可表示,
则 $I$ 是伪滤过的 (\resp 滤过的).
\end{enonce*}

\begin{enonce*}{命题 2.8}\etiquette[2.8]{I.prop2.8}
设 $𝒰$ 是一个宇宙.
$𝒰$-$\Ens$ 中的滤过 $𝒰$-归纳极限与有限投射极限交换.
\end{enonce*}

立即归结为证明滤过 $𝒰$-归纳极限与纤维积交换.
证明留给读者. 可以利用下述引理给出的极限的描述.

\begin{enonce*}{引理 2.8.1}\etiquette[2.8.1]{I.lem2.8.1}
设 $I$ 是一个小滤过范畴, $i\mapsto X_i$ 是从 $I$ 到 $𝒰$-$\Ens$ 的函子.
在不交并 $\amalg_{i\in
 \ob I} X_i$ 上, 设 $R$ 为关系:
\begin{itemize}
\item[(R)] 两个元素 $x_i\in X_i$ 和 $x_j\in X_j$ 有关系,
是指存在对象 $k\in \ob I$ 及两个态射
$u:i\rightarrow k$ 和 $v:j\rightarrow k$,
使得 $x_i$ 和 $x_j$ 分别在转移映射 $u$ 和 $v$ 下在 $X_k$ 中的像相等.
\end{itemize}
则:\pageoriginale
\begin{enumerate}
\renewcommand{\labelenumi}{\theenumi)}
\item $R$ 是等价关系.

\item 商 $\amalg_{i\in \ob I}X_i/R$ 典范同构于 $\varinjlim_I X_i$.

\item 两个元素 $x_i\in X_i$ 和 $x_j\in X_j$
分别在 $R$ 下等价于同一 $X_k$ 中的两个元素.

\item 对每个 $i\in \ob I$, $X_i$ 的两个元素 $\alpha$ 和 $\beta$
在 $R$ 下等价, 当且仅当存在态射
$u:i\rightarrow j$ 使得 $u(\alpha)=u(\beta)$.
\end{enumerate}
\end{enonce*}

断言 1) 来自 (PS 1) \hyperref[I.def2.7]{(2.7)}. 要证明 2), 只须 %% étiquette !
验证 $\amalg_{i\in \ob I}X_i/R$ 具有归纳极限的泛性质
(仅用到 (PS 1)). 断言 3) 来自 $I$ 连通这一事实.
断言 4) 来自 (PS 2).

\begin{enonce*}{推论 2.9}\etiquette[2.9]{I.cor2.9}
设 $%%
\gamma$ 是一种“由有限投射极限定义”的代数结构.
(请读者赋予前一句话以数学含义.
这里只指出, 群、Abel 群、环、模等 %%
结构都属于这一类.) 以 $𝒰$-$\gamma$ 表示
$𝒰$-$\Ens$ 中的 $\gamma$-对象所成的范畴.
将 $𝒰$-$\gamma$ 的每个对象对应于其底集合的函子
与滤过 $𝒰$-归纳极限交换.
因此, $𝒰$-$\gamma$ 中的滤过 $𝒰$-归纳极限与有限投射极限交换.
\end{enonce*}

\begin{enonce*}{推论 2.10}\etiquette[2.10]{I.cor2.10}
$𝒰$-$\Ab$ 中的伪滤过 $𝒰$-归纳极限与有限投射极限交换.
\end{enonce*}

(将指标范畴分解为连通分支, 即归结为滤过极限的情形.)

现在记下\pageoriginale{}一个将不断使用的结果:
设 $C$ 和 $C'$ 是两个 $𝒰$-范畴,
$u$ 和 $v$ 是两个函子 $\xymatrix@C=1.25cm{ C  \ar@<-.6ex>[r]_u & C'
\ar@<-.6ex>[l]_v }$, 且 $u$ 是 $v$ 的\emph{左伴随}.
回忆这意味着, 在取值于 $𝒰$-$\Ens$ 的双函子之间有同构:
$$
\Hom_{C'}(u(X),X')\xrightarrow{\sim}\Hom_{C}(X,v(X')){\quoi}.
$$

\begin{enonce*}{命题 2.11}\etiquette[2.11]{I.prop2.11}
函子 $u$ 与可表示归纳极限交换;
函子 $v$ 与可表示投射极限交换.
\end{enonce*}

这个断言的意思是, 对每个范畴 $I$ 以及每个函子
$G:I\rightarrow C'$ 使得 $G$ 的投射极限可表示, 函子 $v\circ G$ 具有可表示的投射极限,
并且有典范同构
$$
v(\varprojlim G)\longrightarrow \varprojlim(v\circ G){\quoi}.
$$
关于函子 $u$ 的断言由读者自行写出.

最后, 记下指标范畴有积时投射极限的一个计算方法.

\begin{enonce*}{命题 2.12}\etiquette[2.12]{I.prop2.12}
设 $I$ 和 $C$ 是两个范畴, $𝒱$ 是一个宇宙.
假定 $C$ 中的 $𝒱$-投射极限可表示,
而且 $I$ 中存在一族对象 $(i_\alpha)_{\alpha\in A}$,
其中 $A$ 是 $𝒱$ 的元素, 使得:
\begin{enumerate}
\renewcommand{\labelenumi}{\theenumi)}
\item 对每对 $(\alpha,\beta)\in A\times A$, 积 $i_\alpha\times i_\beta$ 可表示,

\item $I$ 的每个对象都存在到至少一个 $i_\alpha$ 的态射.
\end{enumerate}


则对每个从 $I$ 到 $C$ 的反变函子 $F$
$$
F:I^\circ\longrightarrow C{\quoi},
$$
$F$ 的投射极限存在, 并且有一个关于 $F$ 函子性的同构:
$$
\varprojlim F\xrightarrow{\sim} \Ker \left(\prod_{\alpha\in
  A} F(i_\alpha)\rightrightarrows \prod_{(\alpha,\beta)\in A\times A}
F(i_\alpha \times i_\beta)\right){\quoi},
$$
其中\pageoriginale{}两个箭头由积 $i_\alpha\times i_\beta$ 到各因子的投影定义.
\end{enonce*}

% END INCLUDED SGA4_English_current_buildable_source_20260906/source/01/drafts/SGA4_I_LINES_411_824_English.texfrag


% BEGIN INCLUDED SGA4_English_current_buildable_source_20260906/source/01/drafts/SGA4_I_LINES_825_1024_English.texfrag
\section{预层范畴的正合性}\etiquette[3]{I.sec3}

设 $C$ 是一个 $𝒰$-范畴, $\widehat{C}$ 是 $C$ 上的预层范畴.
$\widehat{C}$ 的所有正合性都来自下述命题:

\begin{enonce*}{命题 3.1}\etiquette[3.1]{I.prop3.1}
$\widehat{C}$ 中的 $𝒰$-投射极限与归纳极限都可表示.
对 $C$ 的每个对象 $X$, $\widehat{C}$ 上的函子
$$
F\longmapsto F(X)\qquad F\in \widehat{C}
$$
与归纳极限和投射极限交换.

换句话说, “$\widehat{C}$ 中的归纳极限和投射极限都是逐个自变量计算的”.
我们列举几个推论.
\end{enonce*}

\begin{enonce*}{推论 3.2}\etiquette[3.2]{I.cor3.2}
设 $\gamma$ 是由有限投射极限定义的代数结构.
取值于 $𝒰$-$\gamma$ 的反变函子范畴
等价于 $\widehat{C}$ 中的 $\gamma$-对象范畴.
\end{enonce*}

\begin{enonce*}{推论 3.3}\etiquette[3.3]{I.cor3.3}
$\widehat{C}$ 中同时为单态射和满态射的态射是同构.
$\widehat{C}$ 中的每个态射都唯一分解为一个满态射后接一个单态射.
$\widehat{C}$ 中的可表示归纳极限都是泛的.
$\widehat{C}$ 中的滤过 $𝒰$-归纳极限与有限投射极限交换.
典范函子 $C\rightarrow
\widehat{C}$ 与可表示投射极限交换.
\end{enonce*}

更\pageoriginale{}一般地说, 范畴 $\widehat{C}$
继承了 $𝒰$-集合范畴的一切涉及归纳极限与投射极限的性质.

\setcounter{subsection}{4} \setcounter{subsubsection}{-1}
\subsubsection{}\etiquette[3.4.0]{I.sec3.4.0}
设 $F$ 是 $\widehat{C}$ 的一个对象. 以 $C/F$ 表示如下范畴.
$C/F$ 的对象是由 $C$ 的一个对象 $X$
以及一个从 $X$ 到 $F$ 的态射 $u$ 构成的对.
设 $(X,u)$ 与 $(Y,v)$ 是两个对象.
从 $(X,u)$ 到 $(Y,v)$ 的态射是一个从 $X$ 到 $Y$ 的态射 $g$,
使下图交换:
\[
\xymatrix@=1.25cm{
X \ar[rr]^{g}\ar[rd]_{u}&   & Y\ar[ld]^{v}\\
            & F &.
}
\]

\begin{enonce*}{命题 3.4}\etiquette[3.4]{I.prop3.4}
沿用 \hyperref[I.sec3.4.0]{(3.4.0)} 的记号,
源函子 $C/F\rightarrow
\widehat{C}$ 在 $\widehat{C}$ 中具有可表示的归纳极限.
典范态射
$$
\mathop{\underrightarrow{\text{\rm lim source}}}_{C/F}\;
(.)\longrightarrow F
$$
是同构.
\end{enonce*}


\begin{enonce*}{推论 3.5}\etiquette[3.5]{I.cor3.5}
设 $F$ 和 $H$ 是 $\widehat{C}$ 的两个对象. 有典范同构
$$
\Hom(F,H)\xrightarrow{\sim}\varprojlim_{(X,u)\in \ob(C/F)}
H(X){\quoi}.
$$
\end{enonce*}

\begin{enonce*}[remark]{证明}
此推论立即由 \hyperref[I.prop3.4]{3.4}
和 \hyperref[I.def1.2]{1.2} 得到.
\end{enonce*}

\begin{enonce*}{命题 3.6}\etiquette[3.6]{I.prop3.6}
设\pageoriginale $C$ 是一个 $𝒰$-范畴, $𝒱$ 是包含 $𝒰$ 的宇宙,
$\nobreak{i:\widehat{C}_{𝒰}\rightarrow
\widehat{C}_{𝒱}}$ 是相应预层范畴的自然包含函子
\hyperref[I.condef1.3]{(1.3)}. 函子 $i$ 与归纳极限和投射极限交换.
此外, 对 $\widehat{C}_{𝒰}$ 的每个对象 $F$,
$i(F)$ 的每个子对象都同构于 $F$ 的唯一一个子对象在 $i$ 下的像
(这样就给出了从 $F$ 的子对象集到 $i(F)$ 的子对象集的双射).
\end{enonce*}

\section{筛}\etiquette[4]{I.sec4}

\begin{enonce*}{定义 4.1}\etiquette[4.1]{I.def4.1}
设 $C$ 是一个范畴.
范畴 $C$ 中的\emph{筛（sieve）}是 $C$ 的一个满子范畴 $D$,
具有如下性质: $C$ 中的每个存在到 $D$ 的一个对象的态射的对象都属于 $D$.
如果 $X$ 是 $C$ 的一个对象,
则 $X$ 的筛是指范畴 $C/X$ 中的筛 (这是一种惯常的说法).
\end{enonce*}

设 $𝒰$ 是一个使 $C$ 成为 $𝒰$-范畴的宇宙.
设 $\widehat{C}$ 是相应的预层范畴.
给 $X$ 的每个筛如下地对应 $\widehat{C}$ 中 $X$ 的一个子对象:
给 $C$ 的每个对象 $Y$ 对应由满足对象 $(Y,f)$ 属于此筛的态射
$f:Y\rightarrow X$ 构成的集合.

\begin{enonce*}{命题 4.2}\etiquette[4.2]{I.prop4.2}
上述映射在 %%
$X$ 的筛所成的集合与 $\widehat{C}$ 中 $X$ 的子对象所成的集合之间建立双射.
\end{enonce*}

\begin{enonce*}[remark]{证明}
我们只\sisi{证明下面给出的映射}{说明下面的映射}就是逆映射.
给 $X$ 的每个子函子 $R$ 对应 $C$ 中 $R$ 上的对象所成的范畴 $C/R$
\sisi{(3.4)}{(\hyperref[I.prop3.4]{3.4})}.
立即可验证, $C/R$ 是 $X$ 的筛.
\end{enonce*}

\begin{enonce*}[remark]{注 4.2.1}\etiquette[4.2.1]{I.rem4.2.1}
同样\pageoriginale{}可见, $C$ 的筛与 $C$ 上的“终函子”
($\widehat{C}$ 的“终”对象) 的子函子所成的集合之间有典范双射.
\end{enonce*}

\setcounter{subsection}{2}
\subsection{}\etiquette[4.3]{I.sec4.3}
设 $C$ 是一个 $𝒰$-范畴.
沿用惯常说法, 我们也将范畴 $\widehat{C}$ 中 $X$ 的子对象称为 $X$ 的筛.
由此, 对每个预层 $F$ 及 $X$ 的每个筛 $R$,
可将 $\Hom_{\widehat{C}}(R,F)$ 定义为 %%
从函子 $R$ 到 $F$ 的态射所成的集合.
此外, 有一个关于 $F$ 函子性的典范同构 \hyperref[I.cor3.5]{(3.5)}:
$$
\Hom_{\widehat{C}}(R,F)\xrightarrow{\sim}\varprojlim_{C/R}
F(.){\quoi},
$$
这也给出一个直接定义\nde{更简单地, 将 $C/R$ 与 $R$ 等同, 就有公式 $
\Hom_{\widehat{C}}(R,F)\xrightarrow{\sim}\varprojlim_{R}
F{\quoi}, $ 其中沿用惯常记号, $F$ 表示 $F$ 与典范“源”函子
$R\rightarrow C/X\rightarrow
C$ 的复合.}.
同样, \hyperref[I.prop4.2]{命题 4.2} 使我们可以将函子的通常运算移用于筛.
现列举如下:

\subsubsection{}\etiquette[4.3.1]{I.4.3.1}
\emph{基变换.} 设 $R$ 是 $X$ 的一个筛,
$f:Y\rightarrow X$ 是 $C$ 的对象之间的态射.
纤维积 $\fprod{R}{Y}{X}$ 是 $Y$ 的一个筛,
称为\emph{由 $R$ 经基变换得到的筛}.
$C/Y$ 中相应的子范畴, 就是 $C/X$ 中由 $R$ 定义的子范畴
在 $f$ 所定义的典范函子 $C/Y\rightarrow C/X$ 下的逆像.

\subsubsection{}\etiquette[4.3.2]{I.sec4.3.2}
\emph{序关系、交、并.} $X$ 的子函子之间的包含关系是序关系.
以任意集合为指标的一族筛的并与交,
可以分别定义为相应子预层族的上确界与下确界.

\subsubsection{}\etiquette[4.3.3]{I.sec4.3.3}
\emph{像、生成的筛.}\pageoriginale
设 $(F_\alpha)_{\alpha\in A}$ 是一族预层,
对每个 $\alpha\in A$, 设 $f_\alpha:F_\alpha\rightarrow X$ 是态射,
其中 $X$ 是 $C$ 的一个对象.
这族态射的\emph{像}是各 $f_\alpha$ 的像的并.
因此, 这族态射的像是 $X$ 的一个筛.
特别地, 如果所有 $F_\alpha$ 都是 $C$ 的对象,
就称其像筛为态射 $f_\alpha$ 生成的筛.
范畴 $C/R$ 是 $C/X$ 的满子范畴,
由那些 $X$ 上的对象 $X'\rightarrow X$ 构成,
其中存在从 $X'$ 到某个 $F_{\alpha}$ 的 $X$-态射.

读者可以作为练习, 将这里对子函子定义的关系与运算改用范畴 $C/R$ 的语言表述.
这样便会看到, 这些关系与运算与使 $C$ 成为 $𝒰$-范畴的宇宙 $𝒰$ 无关,
因而对每个范畴都有定义, 无须指明态射集属于哪个宇宙;
这一点其实也可以由 \hyperref[I.prop3.6]{3.6} \textit{预先}看出.

% END INCLUDED SGA4_English_current_buildable_source_20260906/source/01/drafts/SGA4_I_LINES_825_1024_English.texfrag


% BEGIN INCLUDED SGA4_English_current_buildable_source_20260906/source/01/drafts/SGA4_I_LINES_1025_1161_English.texfrag
